\scanpage{107}
% archive foliation: 106
\[
c_{D_0+\omega_1}=c_{D_0}+2D_0(\omega_1)+\omega_1^2
\]
\uline{NB} \inserted{$\omega_1\longmapsto$} $2D_0(\omega_1)+\omega_1^2$ définit un morphisme de faisceaux (ni linéaire, ni ``schématique''\ldots) $\underline\omega\longrightarrow\omega^{\unsure{\otimes2}}$.
Si on remplace $D_0$ par $D'_0=D_0+\bar\omega_1$, $\bar\omega_1\in\Gamma\underline\omega$, \unsure{c'est} augmenté de $\omega_1\longmapsto\omega_1\bar\omega$.

Considérons maintenant une section $\sigma'$ de $P$ sur $X$, donnée \inserted{en termes de $D_0$} par une $\omega_1$ \uline{rationnelle} (pouvant avoir des pôles) --- on se demande s'il $\exists$ une conn. projective rendant $\sigma'$ plate --- laquelle sera unique, car donnée aux pts de régularité de \deleted{$\sigma'$}\inserted{$\omega_1$} par $2D_0(\omega_1)+\omega_1^2$.
Donc une telle conn. projective existe ssi $2D_0(\omega_1)+\omega_1^2$ reste régulière aux pôles de $\omega_1$. Si $D_0=\unsure{d_t|D_t}$ \inserted{$\omega_1=\lambda\,dt$} ça signifie que $2\dfrac{d\lambda}{dt}+\lambda^2$ reste régulière aux pôles de $\lambda$ --- cela implique que le développement polaire de $\lambda$ en un pôle ($t=0$ disons\ldots) est donné par
\[
\lambda=\frac2t+\text{termes d'ordre }\geq1
\]
(le terme const. doit être nul !)
